\documentclass[11pt]{article}
% ===========================================================================
%  Shared preamble for the AP Calculus AB exam-preparation worksheets.
%  Input right after \documentclass. Files beginning with "_" are not compiled
%  on their own. Navy and gold, to match the course's lesson pages.
% ===========================================================================
\usepackage[letterpaper,top=2.3cm,bottom=2.2cm,left=1.7cm,right=1.7cm,headheight=28pt,footskip=30pt]{geometry}
\usepackage{amsmath,amssymb}
\usepackage[T1]{fontenc}
\usepackage[utf8]{inputenc}
\usepackage{lmodern}
\usepackage{xcolor}
\usepackage[most]{tcolorbox}
\usepackage{enumitem}
\usepackage{fancyhdr}
\usepackage{array}
\usepackage{tikz}
\usepackage{pgfplots}
\pgfplotsset{compat=1.18}

\definecolor{navy}{HTML}{161E54}
\definecolor{ink}{HTML}{27357E}
\definecolor{gold}{HTML}{B98A00}
\definecolor{goldbg}{HTML}{FFFDF3}
\definecolor{inkbg}{HTML}{F6F7FD}
\definecolor{rule}{HTML}{C9CFF0}
\definecolor{okgreen}{HTML}{1B7A44}
\definecolor{okbg}{HTML}{F1F8F3}

\pagestyle{fancy}
\fancyhf{}
\fancyhead[L]{\footnotesize Name: \underline{\hspace{4cm}}}
\fancyhead[C]{\textbf{\color{navy}Integration Academy}\\[-2pt]{\scriptsize Calculus AB \textperiodcentered{} AP\textsuperscript{\textregistered} Exam Preparation}}
\fancyfoot[L]{\tiny\color{black!55}\parbox[t]{0.78\textwidth}{AP\textsuperscript{\textregistered} is a registered trademark of the College Board, which was not involved in the production of, and does not endorse, this material. \textcopyright{} Integration Academy.}}
\fancyfoot[R]{\scriptsize\thepage}
\renewcommand{\headrulewidth}{1.2pt}
\renewcommand{\headrule}{\hbox to\headwidth{\color{gold}\leaders\hrule height \headrulewidth\hfill}}
\renewcommand{\footrulewidth}{0pt}
\setlength{\parindent}{0pt}
\setlength{\parskip}{4pt}

% Title block: \wstitle{code}{title}{unit line}
\newcommand{\wstitle}[3]{\begin{center}{\small\bfseries\color{gold}WORKSHEET #1}\\[3pt]{\LARGE\bfseries\color{navy}#2}\\[3pt]{\small\color{black!60}#3}\end{center}\vspace{2pt}}
% Section heading
\newcommand{\sech}[1]{\par\vspace{10pt}{\large\bfseries\color{navy}#1}\par\nobreak\vspace{2pt}{\color{rule}\hrule height0.8pt}\vspace{7pt}}
% A level tag printed beside a group of questions
\newcommand{\level}[1]{\par\vspace{6pt}{\footnotesize\bfseries\color{gold}\MakeUppercase{#1}}\par\nobreak\vspace{2pt}}

\newtcolorbox{ideabox}{colframe=ink,colback=inkbg,coltitle=white,colbacktitle=ink,fonttitle=\bfseries\small,title={Key ideas},arc=2pt,boxrule=0.8pt,left=8pt,right=8pt}
\newtcolorbox{example}[1]{colframe=rule,colback=white,coltitle=navy,colbacktitle=inkbg,fonttitle=\bfseries\small,title={#1},arc=2pt,boxrule=0.8pt,borderline west={3pt}{0pt}{ink},left=10pt,right=8pt}
\newtcolorbox{frbox}[1]{colframe=gold,colback=goldbg,coltitle=white,colbacktitle=gold,fonttitle=\bfseries\small,title={#1},arc=2pt,boxrule=0.8pt,breakable,left=8pt,right=8pt}
\newtcolorbox{solbox}{colframe=okgreen,colback=okbg,coltitle=white,colbacktitle=okgreen,fonttitle=\bfseries\small,title={Worked solutions},arc=2pt,boxrule=0.8pt,breakable,left=8pt,right=8pt}
\newtcolorbox{scorebox}[1]{colframe=navy,colback=white,coltitle=white,colbacktitle=navy,fonttitle=\bfseries\small,title={#1},arc=2pt,boxrule=0.8pt,breakable,left=8pt,right=8pt}

\newcommand{\soln}{\par\smallskip\textbf{\color{ink}Solution.}\ }
\newcommand{\result}[1]{\par\smallskip\textbf{#1}}
% blank working space of a given height
\newcommand{\work}[1]{\par\vspace{#1}}
% a point value in the margin of a scoring guide
\newcommand{\pts}[1]{\hfill{\footnotesize\bfseries\color{gold}[#1]}}

% Coordinate grid for graph questions: \begin{apgraph}{xmin}{xmax}{ymin}{ymax} ... \end{apgraph}
\newenvironment{apgraph}[4]{\begin{center}\begin{tikzpicture}\begin{axis}[
  axis lines=middle, xmin=#1, xmax=#2, ymin=#3, ymax=#4,
  grid=both, grid style={black!12}, xtick distance=1, ytick distance=1,
  ticklabel style={font=\scriptsize}, xlabel={$x$}, ylabel={$y$},
  width=10cm, height=7.2cm, samples=80, clip=false,
  every axis plot/.append style={line width=1.1pt, color=ink}]}{\end{axis}\end{tikzpicture}\end{center}}
\newcommand{\closeddot}[2]{\addplot[only marks, mark=*, mark size=2.2pt, color=ink] coordinates {(#1,#2)};}
\newcommand{\opendot}[2]{\addplot[only marks, mark=*, mark size=2.2pt, color=ink, mark options={fill=white}] coordinates {(#1,#2)};}

% Multiple-choice options: four across, or stacked when the options are long.
\newcommand{\opts}[4]{\par\vspace{3pt}\begin{tabular}{@{}p{0.235\linewidth}p{0.235\linewidth}p{0.235\linewidth}p{0.235\linewidth}@{}}(A)\ #1 & (B)\ #2 & (C)\ #3 & (D)\ #4\end{tabular}\par\vspace{9pt}}
\newcommand{\optsv}[4]{\par\vspace{3pt}\begin{tabular}{@{}l@{\ \ }p{0.9\linewidth}@{}}(A) & #1\\(B) & #2\\(C) & #3\\(D) & #4\end{tabular}\par\vspace{9pt}}

\fancyhead[R]{\footnotesize\color{black!60}Unit 1 test: answer key}
\begin{document}

\wstitle{1.11}{Unit 1 Test: Answer Key}{Unit 1: Limits and Continuity \textperiodcentered{} 20 points}

\begin{scorebox}{Answers}
\begin{center}\renewcommand{\arraystretch}{1.3}
\begin{tabular}{|c|c|c|c|c|c|c|c|c|c|c|}\hline
Question & 1 & 2 & 3 & 4 & 5 & 6 & 7 & 8 & 9 & 10\\\hline
Answer & A & B & D & B & B & A & C & A & C & A\\\hline
Question & 11 & 12 & 13 & 14 & 15 & 16 & 17 & 18 & 19 & 20\\\hline
Answer & D & B & C & D & A & C & B & D & A & C\\\hline
\end{tabular}\end{center}
\end{scorebox}

\begin{solbox}
\begin{enumerate}[leftmargin=*,itemsep=3pt]
  \item $\dfrac{(x-2)(x+2)}{(x-2)(x+3)}=\dfrac{x+2}{x+3}\to\dfrac45$. \textbf{(A)} \hfill{\footnotesize Lesson 1.5}
  \item Both branches approach the open circle at height 0; the filled dot at height 2 is $q(-2)$. \textbf{(B)} \hfill{\footnotesize Lesson 1.2}
  \item The left-hand limit is 3 and the right-hand limit is 4, so the limit does not exist. \textbf{(D)} \hfill{\footnotesize Lesson 1.2}
  \item $\dfrac34\cdot\dfrac{\sin(3x)}{3x}\to\dfrac34$. \textbf{(B)} \hfill{\footnotesize Lesson 1.6}
  \item Equal degrees: $\dfrac62=3$. \textbf{(B)} \hfill{\footnotesize Lesson 1.9}
  \item $\sqrt{16x^2+3}=-x\sqrt{16+\tfrac{3}{x^2}}$ for $x<0$, so the limit is $\dfrac{-4}{2}=-2$. \textbf{(A)} \hfill{\footnotesize Lesson 1.9}
  \item $1+k=3+2k$ gives $k=-2$. \textbf{(C)} \hfill{\footnotesize Lesson 1.8}
  \item $\dfrac{(x-5)(x+5)}{(x-5)(x+1)}$: the factor cancels, and the limit $\tfrac{10}{6}$ exists. \textbf{(A)} \hfill{\footnotesize Lesson 1.7}
  \item The sign changes on $[0,2]$, $[2,5]$ and $[5,8]$: at least 3 zeros, by the Intermediate Value Theorem. \textbf{(C)} \hfill{\footnotesize Lesson 1.10}
  \item $\dfrac{\sqrt x-3}{(\sqrt x-3)(\sqrt x+3)}\to\dfrac16$. \textbf{(A)} \hfill{\footnotesize Lesson 1.5}
  \item $\dfrac{3(2)-(-1)^2}{2+(-1)}=\dfrac51=5$. \textbf{(D)} \hfill{\footnotesize Lesson 1.4}
  \item Both bounds tend to 4, so by the squeeze theorem the limit is 4. \textbf{(B)} \hfill{\footnotesize Lesson 1.6}
  \item Numerator near 4; denominator small and negative. \textbf{(C)} \hfill{\footnotesize Lesson 1.9}
  \item As $x\to1^+$, $q(x)=5-x\to4$ through values less than 4, and $\displaystyle\lim_{u\to 4^-}q(u)=5-4=1$. \textbf{(D)} \hfill{\footnotesize Lesson 1.2}
  \item $\dfrac{x-1}{(x-1)(x+1)}=\dfrac{1}{x+1}$ for $x\neq1$: a hole at $x=1$, an asymptote at $x=-1$. \textbf{(A)} \hfill{\footnotesize Lesson 1.7}
  \item A table at $x=\pm0.001$ gives $2.002$ and $1.998$. \textbf{(C)} \hfill{\footnotesize Lesson 1.3}
  \item A table at $x=\pm0.01$ gives $4.4997$. \textbf{(B)} \hfill{\footnotesize Lesson 1.3}
  \item The entries on both sides approach 5; $g(2)=1$ is the function value. \textbf{(D)} \hfill{\footnotesize Lesson 1.3}
  \item A table at $x=1\pm0.001$ gives $0.9995$ and $1.0005$. \textbf{(A)} \hfill{\footnotesize Lesson 1.3}
  \item $h$ is continuous and $h(1)<0<h(4)$, so the Intermediate Value Theorem gives a zero. It does not give uniqueness. \textbf{(C)} \hfill{\footnotesize Lesson 1.10}
\end{enumerate}
\end{solbox}

\begin{scorebox}{Using your score}
One point per question. 17 or more: ready to move on. 13 to 16: rework the lessons named beside the questions you missed, then try the online practice for those lessons. 12 or fewer: review the unit summary pages before continuing to Unit 2.
\end{scorebox}
\end{document}
